\documentclass[11pt]{article}
\usepackage[margin=1in]{geometry}
\usepackage{iftex}
\ifXeTeX
  \usepackage{fontspec}
  \setmainfont{lmroman10-regular.otf}[BoldFont=lmroman10-bold.otf,
    ItalicFont=lmroman10-italic.otf,BoldItalicFont=lmroman10-bolditalic.otf]
  \newcommand{\vn}[1]{#1}
\else
  % pdfLaTeX (arXiv): T5 for the Vietnamese author name, T1 elsewhere.
  \usepackage[T5,T1]{fontenc}
  \usepackage[utf8]{inputenc}
  \usepackage{lmodern}
  \newcommand{\vn}[1]{{\fontencoding{T5}\selectfont #1}}
\fi
\usepackage{amsmath,amssymb,amsthm,booktabs,array,float}
\usepackage[hidelinks]{hyperref}
\newtheorem{theorem}{Theorem}
\newtheorem{lemma}{Lemma}
\newtheorem{corollary}{Corollary}
\newtheorem{proposition}{Proposition}
\theoremstyle{remark}
\newtheorem{remark}{Remark}
\newcommand{\ind}{\mathbf{1}}
\newcommand{\INF}{\infty}
\newcommand{\rfull}{\rho_{\rm full}}
\newcommand{\rfix}{\rho_{\rm fixed}}
\setlength{\emergencystretch}{2em}
\title{Schema Width Bounds the Benefit of Editable Keys\\
for Plurality-Based Functional-Dependency Checks}
\author{\vn{Trần Vũ Bảo Minh}\\[0.4em]\small Independent Researcher}
\date{November 7, 2026}
\begin{document}
\maketitle
\begin{abstract}
A plurality-based functional-dependency check accepts a target's dependent value
when it is the unique plurality among rows that agree with the target on a
determinant. We study the minimum number of context-cell edits that makes every
listed check accept a protected, initially rejected target, and ask how much
cheaper this becomes when determinant (key) cells may change as well as
dependent cells. For a family $F$ of determinants sharing one dependent
attribute, let $w(F)$ be the largest inclusion antichain among the determinants
that contain a common key. We prove that the fixed-key optimum is at most
$w(F)$ times the editable-key optimum, and that every family has a binary table
attaining exactly this factor. Hence the two optima agree on every such table if
and only if $F$ is laminar. The upper bound uses a margin-capped exchange that
replaces each key edit by at most one dependent-cell flip per inclusion chain.
The classification concerns the listed checks, not logical FD equivalence, and
laminarity does not make optimization easy: acceptance within a budget remains
NP-complete for singleton determinants. For the reverse direction we give an
exact formula. Exhaustive finite checks cover 81,893 rejected table/family
pairs. In a reused-source study on Beers, Flights and Rayyan all nine check
families are laminar, which explains, post hoc, why a predicted accepting gain
from editable keys could not occur.
\end{abstract}

\section{Introduction}
Functional dependencies (FDs) describe agreement between attributes. A simple
detector judges an observed dependent value by comparing it with the unique
plurality among rows that share its determinant values. Such a decision depends
on context. Editing another row's dependent cell changes a vote; editing one of
its key cells moves the row into or out of the target's group. An audit that
trusts all key cells therefore answers a different question from one that
allows them to be wrong.

We measure this as an exact decision-sensitivity problem. The target row is
immutable, every changed context cell costs one, and the edited table need not
satisfy any FD. Let $\rfix$ be the minimum cost of making all checks accept a
rejected target when only dependent cells may change, and $\rfull$ the minimum
when key cells may change too. Always $\rfull\le\rfix$. The question is how
large the gap can be, and which schemas allow a gap at all.

\paragraph{A four-row example.}
Take rules $(a,b)\to y$ and $(a,c)\to y$ and the binary table in
Table~\ref{tab:example}. Each target group contains the target and one
rival, so both checks tie and reject. With fixed keys, $r_1.y$ and $r_2.y$
must both change: $\rfix=2$. Changing only $r_3.a$ to $0$ moves $r_3$
into both groups at once, giving $\rfull=1$. The two determinants share the key
$a$ but neither contains the other. Our main result shows that this is the
only source of a gap, and measures exactly how large it can become.

\begin{table}[H]
\centering\small
\begin{tabular}{lrrrr}
\toprule
Row & $a$ & $b$ & $c$ & $y$\\
\midrule
Target (protected) &0&0&0&0\\
$r_1$ &0&0&1&1\\
$r_2$ &0&1&0&1\\
$r_3$ &1&0&0&0\\
\bottomrule
\end{tabular}
\caption{Constructed illustration, not source data. One key edit replaces two
dependent-cell edits.}\label{tab:example}
\end{table}

\paragraph{Contributions.}
\begin{enumerate}
\item For every listed family $F$ and every initially rejected target we prove
$\rfull\le\rfix\le w(F)\,\rfull$, where $w(F)$ is the maximum, over keys
$j$, of the inclusion width of the determinants containing $j$
(Theorem~\ref{thm:exchange}). The proof converts any optimal editable-key
witness into a fixed-key witness.
\item The factor is sharp for each family separately: some binary table has
$\rfull=1$ and $\rfix=w(F)$ (Theorem~\ref{thm:sharp}). Consequently the two
optima coincide on every table with a rejected target exactly when $F$ is laminar
(Corollary~\ref{cor:laminar}).
\item The equality is a statement about worst cases and edit spaces, not about
tractability. Making a rejected target accept within a budget is NP-complete
for binary singleton schemas, which are laminar (Theorem~\ref{thm:hardness}).
For the opposite direction, rejection, an exact formula shows that keys help
only through singleton groups (Theorem~\ref{thm:reject}).
\item Locked exhaustive checks of raw edit spaces agree with every finite
prediction. A post hoc reanalysis of an earlier source study shows that its
schemas were all laminar, so a predicted accepting gain was impossible; we
report this design error rather than changing the rules after the fact.
\end{enumerate}

The chain-partition step uses Dilworth's theorem~\cite{dilworth}; nested-set
dominance also appears in laminar election-manipulation
algorithms~\cite{emmi2026}. We claim the specific exchange ratio, its sharpness
for each family, and the laminar characterization, not these techniques.
A targeted search found no identical statement (Section~\ref{sec:related});
this is not a historical-priority guarantee. A radius certifies stability of a
decision under a stated edit model. It does not say which cells are wrong and
is not evidence of better cleaning accuracy.

\section{Model}
Let $T=(t,r_1,\ldots,r_{n-1})$ be identified rows with $q\ge1$ categorical key
attributes and one dependent attribute $y$. Duplicate rows are allowed. The RHS
domain has $d\ge2$ values; each key domain contains the target value and an
alternative. By relabeling each column, assume the target is all zero. A
nonempty family $F$ of distinct nonempty subsets of $\{1,\ldots,q\}$ lists the
checks $X\to y$, $X\in F$. Define
\[
 G_X(T)=\{r\in T:r[X]=t[X]\},\qquad c_{Xv}(T)=|\{r\in G_X(T):r_y=v\}|.
\]
The target is accepted iff $c_{X0}(T)>c_{Xv}(T)$ for every $X\in F$ and
$v\ne0$. Ties reject, and a plurality need not be a majority when $d>2$.

An edited table keeps the same rows and leaves $t$ unchanged. Context cells may
independently take any value in their domains. The cost is the cell Hamming
distance $\delta(T,T')$. For a target whose decision is to be flipped, $\rfull$ is
the minimum cost of a flip and $\rfix$ the minimum over edits that keep every
key cell; an impossible flip has cost $\INF$. There are no insertions,
deletions, unequal costs, validity constraints or requirements on other
targets. These restrictions are used in the proofs.

For the accepting direction write $M(r)=\{j:r_j=0\}$ for the zero-key mask of a
context row and define its contribution to rule $X$ and rival $v\ne0$ as
\[
 a_{Xv}(r)=\ind[X\subseteq M(r)]\bigl(\ind[r_y=0]-\ind[r_y=v]\bigr).
\]
The margin is $m_{Xv}(T)=1+\sum_{r\ne t}a_{Xv}(r)$, where the $1$ is the
protected target's vote. Acceptance means $m_{Xv}\ge1$ for all $X,v$. For an
initially rejected target both radii are finite: changing every nonzero context
$y$ to $0$ accepts.

\section{Schema width and the accepting gap}\label{sec:width}
Throughout this section the target is initially rejected, and $\rfull$, $\rfix$
are the costs of making it accept.

For a key $j$ let $F_j=\{X\in F:j\in X\}$, partially ordered by inclusion.
Let $w_j$ be the size of a largest antichain in $F_j$ (zero if $F_j$ is
empty) and
\[
 w(F)=\max_j w_j\ \ge 1.
\]
$F$ is \emph{laminar} if any two intersecting determinants are nested.
Table~\ref{tab:widths} lists simple cases. The last row uses Sperner's theorem
on the subsets of the remaining $q-1$ keys.

\begin{table}[H]
\centering
\begin{tabular}{lcc}
\toprule
Listed family $F$ & $w(F)$ & Laminar\\
\midrule
Singletons $\{1\},\ldots,\{q\}$ & 1 & yes\\
A chain $\{1\}\subset\{1,2\}\subset\cdots\subset\{1,\ldots,q\}$ & 1 & yes\\
$\{a,b\},\{a,c\}$ & 2 & no\\
All $\binom{q}{2}$ pairs, $q\ge3$ & $q-1$ & no\\
All nonempty subsets & $\binom{q-1}{\lfloor (q-1)/2\rfloor}$ & no for $q\ge3$\\
\bottomrule
\end{tabular}
\caption{Width of some determinant families.}\label{tab:widths}
\end{table}

\begin{lemma}[Normalization]\label{lem:norm}
For every original context row $r$ and destination $u$ there is a destination
$c$ with $\delta(r,c)\le \delta(r,u)$ and $a_{Xv}(c)\ge a_{Xv}(u)$ for all $X,v$, of
one of two forms: (i) $c=r$ and $r_y\ne0$; or (ii) $c_y=0$,
$M(c)\supseteq M(r)$, and every changed key of $c$ is set to $0$.
\end{lemma}
\begin{proof}
Suppose $r_y\ne0$ and $u_y=r_y$. If $u=r$, take form (i). Otherwise $u$
changes a key, so $\delta(r,u)\ge1$; take $c$ with the original keys and $c_y=0$,
which costs one. Every contribution of $c$ is nonnegative and every
contribution of $u$ nonpositive.

In all other cases let $c_y=0$, set the keys in $M(r)\cup M(u)$ to $0$, and
keep the original value elsewhere. The key cells changed by $c$ lie in
$M(u)\setminus M(r)$ and are also changed by $u$. If $r_y=0$ the $y$ cost of
$c$ is zero; otherwise it is one and $u_y\ne r_y$ also costs one. If $u_y=0$,
every group containing $u$ contains $c$. If $u_y\ne0$, the contributions of
$u$ are nonpositive while those of $c$ are nonnegative. Since margins are sums over rows, rows can be replaced
independently without increasing cost or decreasing any margin.
\end{proof}

Thus some optimal editable-key witness uses only two row actions: keep a rival
row, or set its $y$ to $0$ and add zero-key matches. A key edit never removes a
rival from a group. Each row has at most $1+2^{q-|M(r)|}$ actions.

\begin{theorem}[Exchange bound]\label{thm:exchange}
For every listed family and every table whose target is initially rejected,
\[
 \rfull\le\rfix\le w(F)\,\rfull.
\]
More precisely, a normalized accepting witness with $R$ dependent-cell edits
and $K_j$ edits in key column $j$ yields a fixed-key accepting witness of cost
at most $R+\sum_j w_jK_j$.
\end{theorem}
\begin{proof}
The first inequality holds because fixed-key edits are a subset. For the second
take an optimal editable-key witness, normalized by Lemma~\ref{lem:norm}. By
Dilworth's theorem~\cite[Theorem 1.1]{dilworth}, partition each $F_j$ into $w_j$
inclusion chains.

Apply the $R$ dependent-cell edits first and copy the result into two tables:
$H$, which will execute the key edits, and $L$, which keeps the original keys
and may flip further nonzero $y$ values to $0$. Maintain for all $X,v$
\begin{equation}\label{eq:inv}
 \min\bigl(m_{Xv}(L),1\bigr)\ \ge\ \min\bigl(m_{Xv}(H),1\bigr).
\end{equation}
Initially equality holds. Consider the next key edit in $H$, in column $j$. By
normalization it sets a key to $0$ in a row whose $y$ is already $0$. The row
joins exactly the groups of some determinants in $F_j$, the \emph{affected}
rules, and each of their margins rises by one. No other margin changes.

For each chain of $F_j$, look at its affected rules whose $L$-group still
contains a context row with nonzero $y$. If there is none, do nothing.
Otherwise let $X$ be the largest such determinant and flip one nonzero-$y$ row
of $G_X(L)$ to $0$. For any other such rule $Y$ in the chain, $Y\subseteq X$,
so $G_Y(L)\supseteq G_X(L)$ contains the flipped row; each of its margins
rises by at least one (by two for the row's former value). An affected rule
whose $L$-group has no nonzero-$y$ context row already has $m_{Xv}(L)\ge1$ from
the target. A flip to $0$ never lowers any margin.

Hence for each affected margin either $m_{Xv}(L)\ge1$ already, or
$m_{Xv}(L)\ge m_{Xv}(H)$ by~\eqref{eq:inv} and $L$ gains at least as much as
$H$. Unaffected margins of $H$ do not change. So~\eqref{eq:inv} persists, using
at most $w_j$ flips for this key edit. Flipped rows had nonzero $y$, so they are
distinct from each other and from the $R$ edited rows. When all key edits are
processed, $H$ accepts; by~\eqref{eq:inv} so does $L$, which never changed a key
or the target. Its cost is at most
$R+\sum_jw_jK_j\le w(F)(R+\sum_jK_j)=w(F)\,\rfull$.
\end{proof}

The cap at one in~\eqref{eq:inv} is needed: a group with no rival left in $L$
may have a smaller margin than in $H$ yet already accepts.

\begin{theorem}[Sharpness for every family]\label{thm:sharp}
For every listed family $F$ there is a binary table with an initially rejected
target such that $\rfull=1$ and $\rfix=w(F)$.
\end{theorem}
\begin{proof}
Choose $j$ with $w_j=w=w(F)$ and an antichain $A_1,\ldots,A_w$ in $F_j$.
In the rows below, every key outside the stated zero mask is $1$.
\begin{itemize}
\item For each $i$, a row $r_i$ with $y=1$ and zero mask exactly $A_i$.
\item A row $p$ with $y=0$ and every key $0$ except key $j$.
\item For each $C\in F$ that contains none of the $A_i$, $w$ support rows with
$y=0$ and zero mask exactly $C$.
\end{itemize}
The group of $A_i$ is $\{t,r_i\}$: $r_l$ with $l\ne i$ would need
$A_i\subseteq A_l$; $p$ fails key $j\in A_i$; a support row for $C$ would need
$A_i\subseteq C$. So each such check ties and rejects.

Every other check accepts. If $C$ contains no $A_i$, its group has the target
and $w$ support rows with $y=0$, against at most $w$ rivals. If
$A_i\subseteq C\ne A_i$ and a rival $r_l$ were in $G_C$, then
$A_i\subseteq C\subseteq A_l$, forcing $l=i$ and $C=A_i$; so $G_C$ has no rival.

Changing $p_j$ to $0$ puts $p$ into every group, making every critical group
$2{:}1$ and harming no check; as the target was rejected, $\rfull=1$. With
fixed keys each critical group contains only the target and its own $r_i$, so
the $w$ distinct cells $r_i.y$ must all change to $0$ (for $d>2$ a third value
still ties), and changing them suffices.
Thus $\rfix=w$. The table has at most $w+1+w|F|$ context rows.
\end{proof}

\begin{corollary}[Laminar characterization]\label{cor:laminar}
The following are equivalent: (1) $\rfix=\rfull$ for every table with an initially rejected target;
(2) $w(F)=1$; (3) $F$ is laminar.
\end{corollary}
\begin{proof}
Two intersecting incomparable determinants share some key $j$ and form a
two-element antichain in $F_j$, and conversely. This gives (2)$\Leftrightarrow$(3).
Theorems~\ref{thm:exchange} and~\ref{thm:sharp} give (1)$\Leftrightarrow$(2).
\end{proof}

The restriction to rejected targets is necessary. For the single rule $x\to y$,
target $(0,0)$ and one context row $(1,1)$, the target is accepted; one key
edit makes it reject, but with fixed keys its group stays $\{t\}$ and
$\rfix=\INF$ (Theorem~\ref{thm:reject}), although $\{x\}$ is laminar.

\begin{remark}[What the result does and does not say]
The factor $w(F)$ is the largest \emph{possible} gain; nonlaminar schemas have
many tables with no gain. For laminar schemas the editable-key search can be
replaced by the fixed-key search, whose only variables are context dependent
cells. That reduction removes key variables but not hardness
(Theorem~\ref{thm:hardness}). For general $F$, any exact fixed-key value is a
$w(F)$-approximation of $\rfull$. The exchange needs an optimal editable-key
witness as input, so it is not itself an approximation algorithm.
\end{remark}

\begin{remark}[Listed checks versus logical FDs]
The result concerns the listed plurality checks, not the FDs they express:
logically equivalent families can behave differently. Add rule $\{a\}$ to
Table~\ref{tab:example} together with two rows $(a,b,c,y)=(0,1,1,0)$. The
$\{a\}$ group has three zeros and two ones and accepts, while the two original
groups still tie, so $F=\{\{a\},\{a,b\},\{a,c\}\}$ still has $\rfix=2$,
$\rfull=1$ and $w(F)=2$. As FDs, $a\to y$ implies $ab\to y$ and $ac\to y$, so
$F$ is equivalent to $\{\{a\}\}$; yet the single check $\{a\}$ has width one
and already accepts this table.
\end{remark}

\section{Complementary results}
\subsection{Making an accepted target reject}
Assume the target is accepted. Write $\alpha_X=c_{X0}(T)$,
$\beta_X=\max_{v\ne0}c_{Xv}(T)$ and $h_X(r)=|\{j\in X:r_j\ne0\}|$.

\begin{theorem}[Exact rejection radius]\label{thm:reject}
Let $\kappa_X=\lceil(\alpha_X-\beta_X)/2\rceil$ if $|G_X(T)|\ge2$, and
$\kappa_X=\min_{r\ne t}\bigl(h_X(r)+\ind[r_y=0]\bigr)$ if $G_X(T)=\{t\}$, with an
empty minimum equal to $\INF$. Then $\rfull=\min_X\kappa_X$. With fixed keys,
$\rfix$ is the minimum over nonsingleton groups only.
\end{theorem}
\begin{proof}
In a nonsingleton group a key edit changes a target/rival margin by at most one
and an RHS edit by at most two, so rejection costs at least
$\lceil(\alpha_X-\beta_X)/2\rceil$. Changing that many zero-valued context RHS cells in
the group to a strongest rival attains it. There are enough, because
acceptance with a nonsingleton group gives $\alpha_X\ge2$ and
$\lceil(\alpha_X-\beta_X)/2\rceil\le \alpha_X-1$. A singleton group must receive a context
row with nonzero $y$, which pays for its $h_X(r)$ nonmatching keys and, if
needed, its RHS; editing just that row creates a tie. Global rejection is a
disjunction over rules, so the minimum over rules is exact. With fixed keys a
singleton group cannot change.
\end{proof}

Thus editable keys strictly lower the rejection radius iff the cheapest
singleton import costs less than every nonsingleton margin cost. The two
decision directions behave differently: rejection depends on singleton groups,
acceptance on the width of the family.

\subsection{Complexity of acceptance}
\begin{theorem}[Binary singleton acceptance is NP-complete]\label{thm:hardness}
Deciding whether a rejected target can be accepted within budget $B$ is
NP-complete in combined complexity, even with binary columns, singleton
determinants and initially tied target groups, with fixed or editable keys.
\end{theorem}
\begin{proof}
Reduce from Set Cover~\cite{karp} with universe $U=\{1,\ldots,q\}$ and sets
$S_1,\ldots,S_\ell$ covering $U$. Let $c_g=|\{i:g\in S_i\}|$ and $\tau$ the
minimum cover size. Use binary keys $x_1,\ldots,x_q$, rules $x_g\to y$ and an
all-zero target. For each $S_i$ add a row with $y=1$ and $x_g=0$ exactly for
$g\in S_i$; for each $g$ add $c_g-1$ rows with $y=0$, $x_g=0$ and all other
keys $1$. Every group has $c_g$ zeros and $c_g$ ones. Flipping the $y$ cells of
a cover's rows accepts, so both optima are at most $\tau$.

Conversely, take an accepting table of cost $K$. For each changed key cell in
column $x_g$ choose an input set containing $g$; for each set row whose final
$y$ is $0$ choose its set. At most $K$ sets are chosen. If column $x_g$ has no
key edit, its group membership is unchanged and, starting from a tie, some
original $y=1$ member, a set row containing $g$, must end at $0$. So the
chosen sets cover $U$ and $\tau\le K$. Both optima equal $\tau$. A destination
table is a polynomial certificate.
\end{proof}

Singleton families are laminar, consistent with both optima being $\tau$. For
fixed $q$ and $d$ a standard dynamic program over the $|F|(d-1)$ margin
coordinates, using the actions of Lemma~\ref{lem:norm}, is polynomial in $n$.
We state this as an elementary observation and have not implemented it.

\subsection{Decision distance is not repair distance}
Let $\mathrm{rep}(T)$ be the minimum context-cell cost making the whole table satisfy
every FD, still protecting $t$.
\begin{proposition}\label{prop:separation}
Even for one rule $x\to y$ with binary $y$, both $\rfull/\mathrm{rep}(T)$ and
$\mathrm{rep}(T)/\rfull$ are unbounded over instances where both are positive and finite.
\end{proposition}
\begin{proof}
With target $(0,0)$, $2h$ rows $(0,0)$ and one row $(0,1)$, one edit repairs
the table but Theorem~\ref{thm:reject} gives $\rfull=h$. With target $(0,0)$
and pairs $(i,0),(i,1)$ for $i=1,\ldots,k$, each disjoint conflicting pair needs
its own edit, so $\mathrm{rep}(T)=k$, while moving $(1,1)$ into the singleton target group
creates a tie at cost one.
\end{proof}

\section{Finite checks}
All studies below were run under protocols committed before their
implementations, and implementations committed before execution. Counts reuse
tables and overlap across edit models; they measure coverage of the checks,
not probabilities, and they do not replace the proofs.

\paragraph{E013: width theorem.}
Three grids use an all-zero target, every multiset of context rows of the
given size, and every family of distinct nonempty determinants: binary keys and
RHS with $q\le3$ and up to three context rows; binary keys, ternary RHS,
$q=3$, up to two rows; ternary keys and RHS, $q=2$, up to two rows. Two raw
oracles enumerate all destination tables (editable keys) and all RHS-only
destinations (fixed keys), counting votes directly; neither uses the theorem.
Width is computed by exhaustive antichains and checked against an independent
minimum chain cover. Every optimal editable-key witness is converted by the
procedure of Theorem~\ref{thm:exchange}, checking cost, unchanged keys,
acceptance and the invariant~\eqref{eq:inv} after each key edit.
Table~\ref{tab:e013} gives the totals. All predictions hold; all 62 strict
gains occur in the binary $q=3$, three-row grid. The sharp construction is
checked for every family with $q\le3$ and four larger families with widths up
to six; for these, the lower bound is checked by the isolated-group
certificate rather than raw enumeration.

\begin{table}[H]
\centering\small
\begin{tabular}{lr}
\toprule
Quantity & Count\\
\midrule
Source tables & 1,895\\
Rejected table/family pairs checked & 81,893\\
\quad of which laminar (equality predicted and observed) & 39,379\\
Pairs with $\rfix>\rfull$ & 62\\
Raw destinations, editable keys / fixed keys & 3,892,339 / 14,687\\
Key edits converted / replacement flips & 11,944 / 12,014\\
Single-row normalization checks & 1,641\\
Sharp constructions checked & 139\\
\bottomrule
\end{tabular}
\caption{E013 totals (protocol \texttt{e25655d}, implementation
\texttt{b93695a}).}\label{tab:e013}
\end{table}

\paragraph{E012: hardness reduction.}
Every coverable multiset of one to three subsets of universes of size one to
three (111 instances) is checked by raw binary enumeration: all 222 radii and
147,523 destinations across both edit models agree with independent minimum
covers, and a cover is extracted from all 25,490 accepting destinations.

\paragraph{E008--E010: earlier formulas and formulations.}
E008 compares canonical-action programs with unrestricted raw edits on 5,107
tables and six families (30,642 radii agree). E009 checks
Theorem~\ref{thm:reject} on all 18,779 accepted cases of that domain. E010
compares a specialized and a generic MILP on 72 constructed cases with up to
six keys and 128 rows; certified optima agree, and seven generic dense-accepted
cases time out. These timings are two repetitions on one machine and are not
representative.

\section{Reused-source study and its post hoc reinterpretation}
E011 reuses Beers, Flights and Rayyan from the Raha repository~\cite{raha} at
revision \href{https://github.com/BigDaMa/raha/tree/7be1334b8c7bbdac3f47ef514fb3e1e8c5fc181c}{\texttt{7be1334b}}.
The data had been inspected earlier, so this is exploratory reuse, not a blind
holdout. Nine tasks use brewery\_id to check brewery\_name, city and state;
flight to check four time columns; Rayyan's \texttt{jounral\_abbreviation}
(original spelling) to check journal\_issn; and abbreviation and ISSN as two
separate determinants for journal\_title. Every target uses its full dirty
table as context; empty strings are ordinary categories.

There are 18,734 target cells: 13,762 accepted and 4,972 rejected. For accepted
targets all radii were computed. Editable keys turn 1,133 fixed-key infinite
radii finite (Table~\ref{tab:census}) and lower one finite radius, Rayyan
journal\_title row 644, from two to one through a singleton group, as
Theorem~\ref{thm:reject} predicts. Among accepted targets, 1,816 differ from
the reference, and 1,493 of these have radius at least two: a stable
acceptance can still be wrong.

\begin{table}[H]
\centering\small
\begin{tabular}{lrrr}
\toprule
Task & Accepted & Fixed $\INF$ to finite & Finite decrease\\
\midrule
Beers brewery\_name&2410&127&0\\
Beers city&2275&127&0\\
Beers state&2275&127&0\\
Flights sched\_dep\_time&1663&0&0\\
Flights act\_dep\_time&1133&0&0\\
Flights sched\_arr\_time&1411&0&0\\
Flights act\_arr\_time&1166&0&0\\
Rayyan journal\_issn&768&384&0\\
Rayyan journal\_title&661&368&1\\
\midrule
Total&13762&1133&1\\
\bottomrule
\end{tabular}
\caption{Accepted-target census from E011.}\label{tab:census}
\end{table}

For rejected targets, a hash-ordered selection of up to four per task gave 32
targets. Each editable-key and fixed-key MILP ran once with a five-second limit.
All 32 fixed solves and 30 editable-key solves were optimal; every exact pair
had equal cost. The declared prediction of at least one strict accepting
gain in the journal\_title task was therefore not supported. Two editable-key
models, rows 509 and 146 (one-based), timed out without bounds; their fixed
optima are 104.

\paragraph{Post hoc reanalysis.}
After E013, we checked the archived E011 schemas. Every one is a singleton
family, or two disjoint singletons for journal\_title, hence laminar. By
Corollary~\ref{cor:laminar} all 4,972 rejected targets have equal fixed and
editable-key accepting costs; the predicted gain could not occur under this edit
model. Combining equality with the archived certified fixed optima gives
editable-key optimum 104 for rows 509 and 146. This deduction rests on the
proof and archived certificates. The solver records remain timeouts, the source
CSV cache was unavailable, and no solver was rerun. The reanalysis was
written after seeing the results and is labeled exploratory. Testing
accepting gains on real data would require a nonlaminar check family justified
and fixed before outcomes are seen.

\section{Related work}\label{sec:related}
\paragraph{Chains and laminar interventions.}
Dilworth's theorem~\cite{dilworth} equates width with the minimum chain
partition; we use it only for that step. In election manipulation under media
influence, strategies sway voter sets. The slides of Miclăuș, Mureșan and
Popa~\cite{emmi2026} credit Tao et al.~\cite{emmi2021} with a dynamic program
for laminar influence sets, and show prefix dominance on chains of nested
influence sets and chain-decomposition parameterized algorithms. We inspected
the 2026 slides and abstract only; neither full paper was obtained. That work optimizes within
one action space, uses co-winner ties, and its laminar family consists of voter
sets. Here laminarity is over column sets; the induced row groups all contain
the target and so are never disjoint. We compare two edit spaces and give the
exact worst-case ratio for each schema. Gerrymandering~\cite{dey} studies
interventions that move voters between districts, so membership-changing
plurality interventions are established; the inspected problem definitions
aggregate district winners and do not compare edit spaces.

\paragraph{Repair bounds depending on the schema.}
Livshits, Kimelfeld and Roy~\cite{livshits} (Proposition 4.4, Corollary 4.5)
convert between optimal subset and update repairs of FD sets with a factor
equal to the minimum lhs cover $\mathrm{mlc}$, and use it for approximation.
Their objective restores consistency of the whole table; ours flips one
protected decision without requiring consistency
(Proposition~\ref{prop:separation}). The parameters differ: $\{\{a,b\},\{a,c\}\}$
has $\mathrm{mlc}=1$ but $w=2$, and $q$ singletons have $\mathrm{mlc}=q$ but $w=1$.

\paragraph{Causality, vote operations and lobbying.}
Cell-based causality for repairs~\cite{cellcauses}, repair-based causes for
query answers~\cite{bertossi} and view-conditioned causality~\cite{vcc} study
minimum interventions on cells; general vote-operation frameworks~\cite{xia}
optimize constructive and destructive goals. Set Cover incidence constructions
are standard in multi-referendum lobbying~\cite{lobbying}; our hardness proof
credits that technique and adds the argument for arbitrary key edits. FD repair
with trusted cells~\cite{repairs}, certified nearest-neighbor robustness over
repairs~\cite{knn}, and the distinction between dependency strength and
detection value~\cite{wan} are further context; like the latter, our radius is
not an accuracy measure.

\section{Limitations}
The proofs need a common dependent attribute, unit and independent cell costs,
an immutable target with the accepting value, and ties that reject. Different
dependents, weighted costs, validity constraints, insertions or deletions, or
co-winner semantics fall outside the stated results; Theorem~\ref{thm:sharp}
in particular uses ties. The width bound is worst case and says nothing about
how often gaps occur. Laminarity does not give a fast algorithm. The finite
checks are small, the source study reuses inspected data, its rejected sample is
small, and its post hoc reanalysis is not a confirmatory test. All proofs and
checks were produced and reviewed inside one AI-assisted project; no
independent human review has taken place. The literature comparison is
targeted, not systematic.

\section{Reproducibility and provenance}
Commit identifiers refer to the public project repository
\url{https://github.com/minggu151623/invariant-audit}. E013 protocol and implementation commits are \texttt{e25655d} and
\texttt{b93695a}; the proof and the E011 reanalysis script were committed at
\texttt{557f60d} before the script ran. E013 used Python 3.14.5 with the
standard library only. E012 uses \texttt{9f14171}/\texttt{65ad518}, E008
\texttt{efc1365}/\texttt{14fe214}, E009 \texttt{96962be}/\texttt{05a2b07}, E010
\texttt{fbd04cb}/\texttt{2212b58} and E011 \texttt{86dad7d}/\texttt{fd341d4};
E009--E011 used Python 3.13.5, NumPy 2.1.3 and SciPy 1.17.1 (HiGHS). Solver
optima required a checked integer witness and bound agreement within $10^{-7}$;
timeouts stay censored. The supplement contains the E013 protocol, code,
summary, compressed case records, proof and a read-only replay that rechecks
every stored witness and can repeat the frozen raw enumeration. Generative AI was
used substantially for ideation, proofs, code and writing.

\section{Conclusion}
For plurality-based FD checks with a protected target, the benefit of allowing
key edits when accepting a rejected target is governed by one schema parameter:
the largest inclusion antichain among determinants sharing a key. The fixed-key
optimum is within this factor of the editable-key optimum, the factor is
attained for every schema, and the optima always agree exactly for laminar
schemas. This explains a negative empirical outcome and separates the question
of when key errors matter from the question of how hard the optimization is.

\begin{thebibliography}{10}
\raggedright
\bibitem{dilworth} R. P. Dilworth. A Decomposition Theorem for Partially
Ordered Sets. \emph{Annals of Mathematics} 51(1):161--166, 1950.
\href{https://doi.org/10.2307/1969503}{doi:10.2307/1969503}.
\bibitem{livshits} E. Livshits, B. Kimelfeld and S. Roy. Computing Optimal
Repairs for Functional Dependencies. \emph{ACM Transactions on Database
Systems} 45(1), 2020.
\href{https://doi.org/10.1145/3360904}{doi:10.1145/3360904}. Conference
version in PODS 2018. Result numbers cited here follow
\href{https://arxiv.org/abs/1712.07705v1}{arXiv:1712.07705v1}.
\bibitem{emmi2021} L. Tao, L. Chen, L. Xu, S. Xu, Z. Gao, W. Shi and D. Huang.
Hardness and Algorithms for Electoral Manipulation Under Media Influence.
IJTCS-FAW 2021, LNCS 12874, pp. 53--64, 2021.
\bibitem{emmi2026} A. Miclăuș, L. Mureșan and A. Popa. Minimum-Cost
Electoral Manipulation Under Media Influence: Hardness, Approximation, and
Algorithms. SYNASC 2026, presentation slides.
\href{https://easychair.org/smart-slide/slide/TstX}{EasyChair Smart Slide},
accessed November 7, 2026.
\bibitem{dey} P. Dey. Gerrymandering: A Briber's Perspective.
\href{https://arxiv.org/abs/1909.01583v1}{arXiv:1909.01583v1}, 2019.
\bibitem{karp} R. M. Karp. Reducibility Among Combinatorial Problems.
In \emph{Complexity of Computer Computations}, pp. 85--103, 1972.
\bibitem{lobbying} R. Bredereck, J. Chen, S. Hartung, R. Niedermeier,
O. Such\'{y} and S. Kratsch. A Multivariate Complexity Analysis of Lobbying in
Multiple Referenda. AAAI, pp. 1292--1298, 2012.
\href{https://doi.org/10.1609/aaai.v26i1.8248}{doi:10.1609/aaai.v26i1.8248}.
\bibitem{cellcauses} M. Debosschere and F. Geerts.
Cell-Based Causality for Data Repairs. TaPP, 2015.
\bibitem{bertossi} L. Bertossi. Specifying and Computing Causes for Query
Answers in Databases via Database Repairs and Repair Programs.
\emph{Knowledge and Information Systems} 63, 2021.
\href{https://doi.org/10.1007/s10115-020-01516-6}{doi:10.1007/s10115-020-01516-6}.
\bibitem{vcc} A. Meliou, W. Gatterbauer, S. Nath and D. Suciu.
Tracing Data Errors with View-Conditioned Causality. SIGMOD, 2011.
\href{https://doi.org/10.1145/1989323.1989376}{doi:10.1145/1989323.1989376}.
\bibitem{xia} L. Xia. How Many Vote Operations Are Needed to Manipulate a
Voting System? COMSOC 2012.
\href{https://arxiv.org/abs/1204.1231v3}{arXiv:1204.1231v3}.
\bibitem{repairs} G. Beskales, I. F. Ilyas and L. Golab. Sampling the Repairs
of Functional Dependency Violations under Hard Constraints. PVLDB 3(1),
197--207, 2010. \href{https://doi.org/10.14778/1920841.1920870}{doi:10.14778/1920841.1920870}.
\bibitem{knn} A. Z. Fan and P. Koutris. Certifiable Robustness for Nearest
Neighbor Classifiers. ICDT, 2022.
\href{https://doi.org/10.4230/LIPIcs.ICDT.2022.6}{doi:10.4230/LIPIcs.ICDT.2022.6}.
\bibitem{wan} X. Wan and X. Han. Strength Is Not Value: Rethinking Approximate
Functional Dependency Measures for Error Detection.
\href{https://arxiv.org/abs/2610.08189v1}{arXiv:2610.08189v1}, 2026.
\bibitem{raha} M. Mahdavi, Z. Abedjan, R. Castro Fernandez, S. Madden,
M. Ouzzani, M. Stonebraker and N. Tang. Raha: A Configuration-Free Error
Detection System. SIGMOD, 2019.
\href{https://doi.org/10.1145/3299869.3324956}{doi:10.1145/3299869.3324956}.
\end{thebibliography}
\end{document}
